\section{Preliminaries and problem statement}
\label{sec:prelim}

\paragraph{Weighted clique complexes.}
All chain spaces are finite-dimensional over \(\C\).  A vertex-weighted
simplicial complex has positive vertex weights \(w(v)\), and a simplex has
weight \(w(\sigma)=\prod_{v\in\sigma}w(v)\).  We use the normalized
coordinates of King and Kohler~\cite{king2026gapped}, in which the weighted
boundary of an oriented simplex removes a vertex \(v\) with coefficient
\(\pm w(v)\); in matrix form \(\partial_{w,k}=W_{k-1}^{-1}\partial_{1,k}W_k\),
where \(W_k\) is the diagonal matrix of \(k\)-simplex weights.  The weighted
complex is therefore conjugate to the unweighted one and has the same
homology.  We use
\[
\Delta_{X,d}=\partial_{X,d}^*\partial_{X,d}
+\partial_{X,d+1}\partial_{X,d+1}^*,\quad
Z_d(X)=\ker\partial_{X,d},\quad
B_d(X)=\ran\partial_{X,d+1}.
\]
Finite-dimensional Hodge theory~\cite{eckmann1944harmonische,horak2013spectra}
gives the orthogonal decomposition
\(C_d(X)=\cH_d(X)\oplus B_d(X)\oplus\ran\partial_{X,d}^*\) with
\(\cH_d(X)=\ker\Delta_{X,d}\) and \(Z_d(X)=\cH_d(X)\oplus B_d(X)\), and
identifies \(H_d(X)=Z_d(X)/B_d(X)\) with the harmonic space \(\cH_d(X)\).
For a positive semidefinite operator \(A\), write
\(\gamma_+(A)=\min(\spec(A)\setminus\{0\})\) for its smallest positive
eigenvalue, with \(\gamma_+(0)=+\infty\).
All certified boundary matrices and fillings below are rational.  Rank, kernel
dimension, image intersection, and filling validity are unchanged under scalar
extension \(\mathbb Q\subset\mathbb R\subset\mathbb C\), and the real
positive-semidefinite inequalities complexify verbatim.  We state homology
over \(\mathbb C\) while certifying the finite palette over \(\mathbb Q\).

\paragraph{Normalized persistence.}
For \(X_1\subseteq X_2\), define
\begin{equation}
\label{eq:true-ratio}
\beta_d(X_1\to X_2)
=\rank[H_d(X_1)\to H_d(X_2)],\qquad
q_d(X_1,X_2)
=\frac{\beta_d(X_1\to X_2)}{\beta_d(X_1)}.
\end{equation}
Lowe et al.~\cite{lowe2026normalized} write the numerator through harmonic
projectors, as the rank of \(P_{\cH_d(X_2)}\) restricted to \(\cH_d(X_1)\).
The two formulations agree.  For a cycle \(z\in Z_d(X_2)\), the vector
\(P_{\cH_d(X_2)}z\) is the harmonic representative of \([z]\), because
\(z-P_{\cH_d(X_2)}z\in B_d(X_2)\); hence \(z\mapsto P_{\cH_d(X_2)}z\) induces
the isomorphism \(H_d(X_2)\cong\cH_d(X_2)\).  Since
\(\cH_d(X_1)\subseteq Z_d(X_1)\subseteq Z_d(X_2)\) and
\(\cH_d(X_1)\cong H_d(X_1)\), the operator \(P_{\cH_d(X_2)}|_{\cH_d(X_1)}\) is
the induced map \(H_d(X_1)\to H_d(X_2)\) up to these identifications, and its
rank is \(\beta_d(X_1\to X_2)\).

\begin{definition}[Gapped normalized persistence]
\label{def:problem}
An instance of the weighted problem consists of explicit simple graphs
\(G_1\subseteq G_2\) on one labeled vertex set \(V\), one common positive
dyadic weight function \(w:V\to\mathbb Q_{>0}\), a degree \(d\ge1\), positive
rational lower bounds \(\gamma_1,\gamma_2\), and a positive rational accuracy
\(\varepsilon\), all encoded in binary and all at least \(1/\poly(|V|)\).
Put \(X_i=\Cl(G_i)\), using the
weighted-boundary convention above.  The promises are
\[
\beta_d(X_1)>0,
\qquad
\spec(\Delta_{X_i,d})\subseteq\{0\}\cup[\gamma_i,\infty)
\quad(i=1,2).
\]
The task is to output \(\widetilde q\) such that
\(\Pr[|\widetilde q-q_d(X_1,X_2)|\le\varepsilon]\ge2/3\).  Input size is the
total bit length of the graph encodings, weights, \(d\), the two gap bounds,
and \(\varepsilon\).  The unweighted problem is the restriction \(w\equiv1\).
Its fixed decision version distinguishes \(q_d\ge3/4\) from \(q_d\le1/8\).
\end{definition}
This is Problem~9 of Lowe et al.~\cite{lowe2026normalized}, with the same
promises, including their inverse-polynomial lower bounds on the gap and
accuracy parameters, and with the encoding conventions fixed.  The denominator is the
initial Betti number, not the number of simplices, and the statistic counts
exact zero modes rather than eigenvalues below a cutoff.

\paragraph{Register, logical projectors, and gadgets.}
The construction starts from a top-dimensional \emph{register} complex
\(R\), so \(C_{d+1}(R)=0\), and sets \(V=Z_d(R)=\cH_d(R)\).  Concretely,
\(R\) is the clique complex of the join of \(n\) \emph{bowtie} graphs, one
per logical qubit.  A bowtie is the wedge of two \(4\)-cycles; its clique
complex is one-dimensional and connected with \(\beta_1=2\), and the two
normalized oriented petal cycles are the computational basis
\(\ket0,\ket1\).  By the K\"unneth formula for joins, \(R\) has reduced
homology only in degree \(d=2n-1\), and \(V\cong(\C^2)^{\otimes n}\)
isometrically.  Each constraint \(j\) is a rank-one projector \(\pi_j\) on
the cycle space of a constant-size \emph{active} sub-register; after
tensoring with the identity of the unaffected factors it becomes an
orthogonal projector \(\Pi_j\) on \(V\), generally of higher rank.  For a
term set \(A\), put
\begin{equation}
\label{eq:logical-sum}
H_A=\sum_{j\in A}\Pi_j,\qquad
K_A=\ker H_A,\qquad
W_A=\sum_{j\in A}\ran\Pi_j=K_A^\perp ,
\end{equation}
and use a certified logical gap
\begin{equation}
\label{eq:logical-gap}
H_A\succeq g_A(P_V-P_{K_A}),\qquad 0<g_A\le1.
\end{equation}

Each constraint is realized by a weighted clique gadget: a graph containing
the active register, whose additional \emph{private} vertices have weight
\(\lambda\) and are joined to every vertex of the unaffected register
factors, while all register vertices have weight one.  Attaching the gadgets
of all \(j\in A\) to \(R\) gives the clique complex \(X_A\), and \(\Delta_A\)
denotes its full Laplacian on all degree-\(d\) chains.  Thus \(H_A\) acts on
\(V\), while \(\Delta_A\) acts on \(C_d(X_A)\supseteq V\).

\begin{definition}[Independent private interiors]
\label{def:independent}
The gadgets of a term set \(A\) have \emph{independent private interiors}
if private vertices of distinct gadgets are never adjacent.  Every simplex
of \(X_A\) that contains a private vertex then contains private vertices of
exactly one gadget, and
\(C_d(X_A)=C_d(R)\oplus\bigoplus_{j\in A}C^{\rm priv}_{d,j}\) as coordinate
spaces.
\end{definition}
All constructions in this paper satisfy \cref{def:independent}: distinct
gadgets receive disjoint sets of new vertices, and no edge is added between
private vertices of distinct gadgets.

\paragraph{Circuit source.}
For the circuit application, \(J\) inserts the fixed clean string around a
maximally mixed logical input, \(U\) is an exact circuit, and
\[
M=J^*U^*P_{\rm acc}UJ,\qquad S=\ker(I-M).
\]
The initial history kernel will have the full input dimension, while the
final history kernel will be isomorphic to the perfect accepting space
\(S\).
